\documentclass[11pt]{article}
\usepackage[margin=1in]{geometry}
\usepackage{amsmath,amssymb}

\title{Proof of Conjecture 00000000101\\
\large A transitive subshift with exactly the prime least periods and \(1\)}
\author{}
\date{}

\begin{document}
\maketitle

The period of a periodic point means its least positive period. This is
the period of its finite orbit, as distinguished from the multiples that
also fix the point. We construct a two-sided subshift over \(\{0,1\}\).

\section*{The gap rule and transitivity}

For a finite binary word \(w\), record the positive distances between
successive occurrences of \(1\), in order. The resulting list can be
empty. Call \(w\) admissible if each nonempty contiguous block of \(k\)
distances in its list satisfies at least one of these conditions
\begin{itemize}
  \item The distances are all equal to the same prime.
  \item At least one of the distances is at least \(k\).
\end{itemize}
Let \(X\) consist of those \(x\in\{0,1\}^{\mathbb Z}\) whose every finite
contiguous word is admissible. With the product topology, \(X\) is
closed, because a violation is witnessed on finitely many coordinates.
Shifting a sequence preserves this condition, in either direction,
so \(X\) is a subshift. Its language is exactly the admissible words.
Indeed, surround any such finite word with zeros in both directions.
Every complete gap in the resulting bi-infinite point is a gap of
the original finite word, and every gap block is inherited from it.
The resulting point thus belongs to \(X\).

For any admissible words \(u,v\), consider \(u0^m v\), where \(m\) can
be made arbitrarily large. If either word has no \(1\), joining adds
no complete gap. Otherwise let \(k_u,k_v\) be the numbers of complete
gaps in \(u,v\). Choose \(m\ge k_u+k_v\). The new gap from the last
\(1\) in \(u\) to the first \(1\) in \(v\) has length at least
\(m+1\ge k_u+k_v+1\), the total number of complete gaps in
\(u0^m v\). Any contiguous gap block that crosses the joining gap
therefore satisfies the second condition. All other gap blocks were
already admissible. This proves that \(u0^m v\) is admissible for
every sufficiently large \(m\), and hence occurs in \(X\).
For any two nonempty cylinders, take the words specifying them
on their respective coordinate intervals. Place a joined word
realizing the first at its specified interval. The second then occurs
at the specified interval after a nonnegative number of left shifts,
by taking \(m\) sufficiently large. Cylinders form a basis of the
topology, so \(X\) is topologically transitive.

\section*{Least periods}

The all-zero point belongs to \(X\) and has least period \(1\).
For each prime \(p\), the point having ones precisely at the multiples
of \(p\) has every complete gap equal to \(p\), so belongs to \(X\).
Its least period is \(p\), since shifting the \(1\) at \(0\) to a
positive position smaller than \(p\) gives a \(0\).

Conversely, suppose a periodic point of \(X\) contains a \(1\) and
has least period \(n\). Choose a \(1\) at coordinate \(a\) and record
the nonempty list \(G\) of complete gaps from the \(1\) at \(a\) to
the \(1\) at \(a+n\), including the gap that ends at \(a+n\).
The sum of \(G\) is \(n\). By periodicity, \(n+1\) consecutive
copies of \(G\) form a contiguous gap block of this point. That block
has length greater than \(n\) and every distance in it is at most
\(n\). The gap rule forces its entries, and hence those of \(G\),
to be equal to a single prime \(p\). Periodicity extends this list
of gaps in both directions, so the ones occur exactly on
\(a+p\mathbb Z\). This point has least period \(p\), proving that
there are no other least periods.

\section*{Lean verification}

The accompanying Lean~4.19.0 project uses the run-length code of a finite
binary word: either its all-zero length, or its leading and trailing
zero lengths together with its list of complete gaps. The theorem
\verb|transitive_language| proves that the admissible word language
has the zero-padding joining property. This gives transitivity by the
cylinder argument above. The theorem \verb|exact_least_period_spectrum|
checks the spectrum on cyclic gap lists whose every finite repetition
is admissible. A least-period orbit with a \(1\) has a nonempty primitive
cyclic gap list, since a repetition of a shorter list would give a
smaller period equal to that shorter list's sum. Conversely any such
primitive list represents the binary point obtained by repeating its
gaps in both directions. Its least period divides the sum of its gaps,
since the least positive period divides every period by the division
algorithm. If the least period were smaller than the sum, the list
would repeat the shorter gap list of that least period, contradicting
primitivity. The theorem \verb|binary_subshift_least_period_spectrum|
also verifies the exact spectrum directly for integer-indexed binary
points satisfying the gap rule. Run \verb|lake build| in \verb|lean/|
to check the formal proofs and print their axiom dependencies.

\end{document}
